\section{Introduction}
\label{sec:intro}

An overcomplete code places $F$ feature directions in $d<F$ dimensions. Such codes are the object of
study in work on superposition and computation in superposition, and they are what a sparse
autoencoder returns when it decomposes a language model's activations into a dictionary of feature
directions. Whatever else one wants from such a code, a great deal of downstream practice reads
features out of it \emph{linearly}: linear probes, the decoder of the autoencoder itself, steering
vectors, and any diagnostic built from inner products with feature directions.

Linear readout of an overcomplete code cannot be exact. With $F>d$ the directions cannot be
orthogonal, so recovering one feature's coefficient picks up contributions from the others, and the
average squared size of that cross-talk is bounded below by a quantity that depends only on the two
counts,
\[
  W(F,d)\;=\;\frac{F-d}{d(F-1)} .
\]
Reporting how close a code comes to this rank--trace floor is therefore a natural way to ask whether
it is arranged well for linear readout, and ratios of this kind have been read as evidence that
training discovers efficient geometry.

\paragraph{What a reader should take from this paper.} Proximity to the floor is not evidence of
anything a code learned, because it is what an unstructured code gives for free; the quantity that
distance from the floor actually measures is the \emph{dispersion of leverage} across features; and
measured this way, no released sparse-autoencoder dictionary we examined is near the floor --- every
one carries more cross-talk under the best linear readout it admits than a random code of its own
shape does, while an
autoencoder trained by the same recipe on a randomly initialised model is not. The practical
consequence is that the excess is a property of the represented content rather than of the fitting
procedure, and that it moves with depth but not with dictionary width.

\subsection{Why proximity to the floor says less than it appears to}

Our starting point is an identity, not an experiment. For a code $\Phi\in\R^{d\times F}$ read out by
its gain-calibrated pseudoinverse, the ratio of achieved cross-talk to the floor is
\[
  \Rgeom(\Phi)\;=\;\frac{d}{F(F-d)}\left(\sum_{i=1}^{F}\frac{1}{h_i}-F\right),
  \qquad h_i=\phi_i^{\top}(\Phi\Phi^{\top})^{-1}\phi_i ,
\]
and the leverage scores obey $\sum_i h_i=d$ exactly. Cauchy--Schwarz then gives
$\sum_i h_i^{-1}\ge F^2/d$ with equality precisely when every $h_i$ equals $d/F$, so
$\Rgeom\ge1$ with equality exactly when leverage is equalised across features
(Proposition~\ref{prop:leverage}). Distance from the floor is not a separate property that a code
might or might not have alongside its leverage profile: \emph{it is} that profile's dispersion, in a
specific harmonic sense that we make precise and then use.

Two consequences follow immediately and are worth stating before any measurement. First, a code whose
leverage is nearly uniform sits at the floor whether or not it encodes anything, which is why an
i.i.d.\ Gaussian code attains it. Second, the ratio is driven by $\sum_i h_i^{-1}$, a harmonic
quantity, so the spread of leverage can grow substantially without the ratio moving --- a
dissociation we observe and would otherwise have had to report as an anomaly.

\subsection{What we measure, and on what}

We evaluate the diagnostic on two kinds of object, and the division of labour between them is
deliberate.

\emph{Released dictionaries} (Section~\ref{sec:dictionaries}) test whether any of this transfers
beyond a controlled setting. We take the decoder matrices of \SaeDicts{} published sparse
autoencoders --- five model families and three independent training recipes, with $d$ from \SaeDMin{}
to \SaeDMax{} and feature loads from \SaeLoadMin$\times$ to \SaeLoadMax$\times$ --- and compute
$\Rgeom$, the leverage distribution, and the failure threshold from the decoder alone. No model
weights, no activations, no training, and no linear programme are involved, so the whole section
costs seconds per dictionary and can be re-run by a reader.

\emph{Controlled campaigns} (Section~\ref{sec:campaigns}) supply what released artefacts cannot: an
intervention. We train 400 small networks in which the loss, the width, the feature load and the
training sparsity are set by us, so that the effect of each on the resulting code can be measured
rather than inferred, and we retain every trained model so that later analyses need no retraining.

\subsection{Contributions}

\begin{enumerate}[leftmargin=*,itemsep=3pt]

\item \textbf{Distance from the rank--trace floor is leverage dispersion, exactly.} The identity is
elementary and we claim no novelty for the mathematics; what is new is that it does predictive work.
It explains why unstructured codes attain the floor, why trained dictionaries do not, and why
dispersion can grow while the ratio stays flat. We verify it against the implementation on
\SaeIdentityRows{} measured codes, with a maximum deviation of \SaeIdentityGap{}.

\item \textbf{Near-floor geometry is generic, and we measure the baseline rather than assert it.}
\SaeControls{} i.i.d.\ controls, matched in shape and operating sparsity to each dictionary, land at
$\Rgeom=\SaeCtlRgeomMin$ to \SaeCtlRgeomMax{} across loads from \SaeLoadMin$\times$ to
\SaeLoadMax$\times$. Any claim that a learned code is close to optimal for linear readout needs this
comparison, and it is not present in the literature we build on.

\item \textbf{No released dictionary we measured is at the floor.} All \SaeDictsAboveControl{} of
\SaeDicts{} sit above their own matched control, at $\Rgeom=\SaeDictRgeomMin$ to
\SaeDictRgeomMax{}, with leverage dispersion \SaeDictCvMin{} to \SaeDictCvMax{} against the
controls' \SaeCtlCvMin{} to \SaeCtlCvMax{}. Whatever these dictionaries optimise, it is not the
evenness of their own leverage.

\item \textbf{The same holds in sparsity units, which is the form a practitioner can use.} The
diagnostic also returns the sparsity at which some feature is guaranteed to be unrecoverable by any
affine rule. All \SaeFailEarlier{} dictionaries reach it \SaeFailGapMin{} to \SaeFailGapMax{} levels
earlier than the i.i.d.\ code of their own shape and operating sparsity: where a random code is still
safe, a released dictionary already has a feature no affine rule recovers. This is a worst-case
statement about one feature rather than an average over all of them, so its agreement with the ratio
is not automatic --- though both are functionals of the same leverage distribution, and neither is an
external endpoint (Section~\ref{sec:discussion}).

\item \textbf{A matched randomly-initialised control locates the cause.} Autoencoders trained by an
identical recipe on a randomly initialised copy of the same model return to the floor
($\Rgeom=\SaeRandInitRgeomMin$ to \SaeRandInitRgeomMax) in all \SaePairedLayers{} paired layers. The
excess therefore reflects what the model learned, not the act of fitting a sparse autoencoder --- a
distinction an i.i.d.\ control cannot draw, because it does not share the training procedure.

\item \textbf{The loss decides whether a linear interface exists at all.} Across four widths, two
feature loads and two training sparsities, $L^2$ training reaches $\Rgeom=\SALtwoFiftyRgeom$ to
\ScLtwoRgeom{} with its affine recovery frontier collapsed to $1$, while $L^4$ holds
$\SALfourTwoHundredRgeom$ with a frontier that grows from \FullLfourFiftyKappaMin{} to
\ScMeasured{}. This is the largest effect we report and the one that replicates most widely.

\item \textbf{The frontier follows a rate, tested out of sample, and the advantage it buys stops
growing.} Fitted on three widths the frontier implied \ScRegistered{} at $d=400$; we registered that
number in the run configuration before the width was trained and measured \ScMeasured{}. Over the
same four widths the advantage over an untrained code narrows from \MarginFiftyRatio{} to
\MarginFourHundredRatio{}, which we report as a limitation on the preceding claim rather than
alongside it.

\item \textbf{An artefact that can be checked.} The weights of all 400 trained networks are
released, along with the SHA-256 of every third-party dictionary analysed, a decision log recording
what was registered before each campaign, and a register of nine defects found during this work
together with what each invalidated. Section~\ref{sec:limitations} states which readings we
withdrew and why.

\end{enumerate}

\subsection{Scope}

Three limits apply throughout and we state them here rather than at the end. Every statement about
the affine recovery frontier is a worst case over feature amplitudes: it says what no affine rule can
guarantee, not what fails in typical use. The comparison between a network's own decoder and a fitted
affine probe (Section~\ref{sec:probes}) measures what supervised fitting recovers rather than what an
affine map can express, because the probe class contains the network's decoder --- and we verify
that the class membership is not hypothetical: the selection misses the network's own map, on the
very validation objective it maximises, in every $L^4$ model above the tie width. We report the
comparison for completeness and build nothing on it. And the controlled campaigns use one task family at widths up
to $d=400$, so they establish the mechanism, not its magnitude in a deployed model --- which is
precisely why the dictionary measurements are the other half of the paper.
